\renewcommand{\thetheorem}{17.d.\arabic{theorem}}
\setcounter{theorem}{0}

% ---------------------------------------------------------------------------
% Provenance of the prose in this file.
%   % Extractor: <model>   the block below was transcribed from Prof. Biran's
%                          handwritten notes by that model
%   % Creator:   <model>   the block below was written by that model and has no
%                          counterpart in the notes
% A tag holds until the next tag.
% ---------------------------------------------------------------------------

% Extractor: Gemini 3.1 Pro
\section{More on Matrices and Linear Equations}
\label{sec:more_matrices_lineq}

Let $A \in \M_{m \times n}(K)$, and consider the homogeneous system of linear equations $A \cdot x = 0$. We denote the space of solutions of this system by:
\begin{equation*}
    \Sol(A, 0) := \{x \in K^n \mid A x = 0\} = \ker(T_A)
\end{equation*}
Note that by the Rank-Nullity Theorem, \cref{thm:rank_nullity}, we have
\[\dim \Sol(A, 0) = n - \rank(A),\]
because $\dim \ker(T_A) = n - \dim \Img(T_A)$ (and $\dim \Img(T_A) = \rankcol(A) = \rank(A)$ by \cref{rem:rank_notation_and_properties}).

Let $A'$ be a row-reduced echelon matrix that is row-equivalent to $A$ (for example, $A'$ is obtained from $A$ by Gauss elimination). We then know that the number of pivots of $A'$ equals $\rank(A)$ (\cref{rem:rank_notation_and_properties} again), and the number of free variables is $n - \rank(A)$.

We can visualize the structure of $A'$ and its corresponding variables $x_1, \dots, x_n$ as in \cref{fig:rref_pivot_staircase}. The columns containing the leading $1$'s (pivots) correspond to the pivot variables, while the remaining columns correspond to the free variables. Clearly, the solution spaces are identical (this is \cref{thm:ero_preserve_solutions}: row-equivalent systems have the same solutions, applied to the homogeneous system):
\begin{equation*}
    \Sol(A, 0) = \Sol(A', 0)
\end{equation*}

% Creator: Gemini 3.1 Pro , TikZ rendering of the staircase sketch in the notes
\begin{figure}[htbp]
\centering
\begin{tikzpicture}[x=8mm, y=8mm, >=Stealth]
  %, the pivot pattern: column index of the leading 1 in each non-zero row ---
  % rows are drawn top-down; row k has its pivot in column pivotcol{k}
  \def\ncols{9}
  \def\nrows{6}   % 4 non-zero rows + 2 zero rows

  % Shade the pivot columns (2, 3, 6, 8) and the free columns.
  \foreach \c in {2,3,6,8}
    \fill[ThemePurple!12] (\c-1, 0) rectangle (\c, \nrows);
  \foreach \c in {1,4,5,7,9}
    \fill[BurntOrange!10] (\c-1, 0) rectangle (\c, \nrows);

  % The zero rows at the bottom.
  \fill[black!6] (0,0) rectangle (\ncols, 2);

  % The staircase: everything strictly below/left of it is zero.
  \draw[very thick, ThemePurple!80!black]
    (0,\nrows) -- (1,\nrows) -- (1,5) -- (2,5) -- (2,4) -- (5,4) -- (5,3)
    -- (7,3) -- (7,2) -- (\ncols,2);

  % Leading 1's.
  \node at (1.5, 5.5) {$1$};
  \node at (2.5, 4.5) {$1$};
  \node at (5.5, 3.5) {$1$};
  \node at (7.5, 2.5) {$1$};

  % A hint of the free entries to the right of each pivot.
  \node[gray] at (4.0, 5.5) {$*$};   \node[gray] at (6.5, 5.5) {$*$};
  \node[gray] at (8.5, 5.5) {$*$};   \node[gray] at (4.0, 4.5) {$*$};
  \node[gray] at (6.5, 4.5) {$*$};   \node[gray] at (8.5, 4.5) {$*$};
  \node[gray] at (6.5, 3.5) {$*$};   \node[gray] at (8.5, 3.5) {$*$};
  \node[gray] at (8.5, 2.5) {$*$};
  \node[black!45, font=\Large] at (4.5, 1.0) {$0$};

  % Matrix delimiters.
  \draw[thick] (-0.18,0) -- (-0.35,0) -- (-0.35,\nrows) -- (-0.18,\nrows);
  \draw[thick] (\ncols+0.18,0) -- (\ncols+0.35,0) -- (\ncols+0.35,\nrows) -- (\ncols+0.18,\nrows);

  % Column labels: which variable each column belongs to.
  \foreach \c/\lab in {1/{x_1},2/{x_2},3/{x_3},4/{x_4},5/{x_5},6/{x_6},7/{x_7},8/{x_8},9/{x_9}}
    \node[below] at (\c-0.5, 0) {\small $\lab$};

  % Braces classifying the columns.
  \node[below=5mm, ThemePurple!80!black, font=\small\sffamily\bfseries] at (4.5,-0.35)
    {pivot columns: $x_2, x_3, x_6, x_8$ \quad ($r = \rank(A) = 4$ of them)};
  \node[below=9mm, BurntOrange!80!black, font=\small\sffamily\bfseries] at (4.5,-0.35)
    {free columns: $x_1, x_4, x_5, x_7, x_9$ \quad ($\ell = n - \rank(A) = 5$ of them)};

  \node[right=3mm, black!55, font=\small] at (\ncols+0.35, 1.0) {zero rows};
\end{tikzpicture}
\caption{The staircase shape of a row-reduced echelon matrix $A'$, here with
$n = 9$ columns and $\rank(A) = 4$. Every entry below the purple staircase
vanishes. Each of the $r$ non-zero rows contributes exactly one leading $1$,
whose column is a \emph{pivot} column; the remaining $\ell = n - r$ columns are
the \emph{free} columns. Choosing the $\ell$ free variables arbitrarily forces
the $r$ pivot variables, which is precisely the map $\Phi$ below.}
\label{fig:rref_pivot_staircase}
\end{figure}

Denote by $x_{i_1}, \dots, x_{i_\ell}$ the free variables, where $1 \le i_1 < \dots < i_\ell \le n$ and $\ell = n - \rank(A)$. Similarly, denote by $x_{j_1}, \dots, x_{j_r}$ the pivot variables, where $1 \le j_1 < \dots < j_r \le n$ and $r = \rank(A)$.

% Creator: Opus 5
% The map $\Phi$ is defined in running prose, and the proposition immediately
% below is a statement about it, so it wants an environment of its own.
% Unnumbered, to leave Prof. Biran's numbering 17.d.1--17.d.4 intact.
\begin{definition*}[The Free-Variable Parametrisation of the Solution Space]
\label{def:free_variable_parametrisation}
% Extractor: Gemini 3.1 Pro
For each choice of values for the free variables $x_{i_1}, \dots, x_{i_\ell} \in K$, the pivot variables $x_{j_1}, \dots, x_{j_r}$ are uniquely determined. So, we obtain a map $\Phi: K^\ell \to \Sol(A, 0)$, where $\ell = n - \rank(A)$, defined by:
\begin{equation*}
    \Phi(\lambda_1, \dots, \lambda_\ell) := \text{the unique solution of } Ax = 0 \text{ (which is the solution of } A'x = 0 \text{)}
\end{equation*}
with $x_{i_1} = \lambda_1, \dots, x_{i_k} = \lambda_k, \dots, x_{i_\ell} = \lambda_\ell$.
\end{definition*}
% Extractor: Gemini 3.1 Pro

\begin{proposition}[Linearity and Isomorphism of the Solution Map]
% prop:17.d.1
\label{prop:solution_map_isomorphism}
The map $\Phi$ of the Definition on \cpageref{def:free_variable_parametrisation} is linear. Moreover, it is an isomorphism.
\end{proposition}

\begin{proof}
\textbf{Linearity of $\Phi$:} The $k$-th non-zero row of $A'$ gives the following equation:
\begin{equation*}
    x_{j_k} + (\text{linear combination of the free variables that come after } j_k) = 0
\end{equation*}
This allows us to isolate the pivot variable $x_{j_k}$:
\begin{equation*}
    x_{j_k} = - \sum_{\{q \mid 1 \le q \le \ell, \, i_q > j_k\}} a'_{k, i_q} \cdot x_{i_q}
\end{equation*}
This explicitly shows that the pivot variables $x_{j_1}, \dots, x_{j_r}$ depend linearly on the free variables $x_{i_1}, \dots, x_{i_\ell}$. This proves the linearity of $\Phi$.

\textbf{Isomorphism:} To prove that $\Phi$ is an isomorphism, it is enough to show that $\ker(\Phi) = \{0\}$, because $\Phi: K^\ell \to \Sol(A, 0)$ is a linear map between two spaces of the exact same dimension $\ell$ (the target has dimension $\ell = n - \rank(A)$ by the count at the start of this section; for maps between spaces of equal finite dimension, injectivity already gives bijectivity by \cref{cor:equal_dim_iso} \textbf{(c)}, and a bijective linear map is an isomorphism by \cref{lem:bijective_is_iso}).

But this clearly holds: if $\lambda = (\lambda_1, \dots, \lambda_\ell) \in \ker(\Phi)$, which means $\Phi(\lambda_1, \dots, \lambda_\ell) = 0$, then by the definition of $\Phi$, we must have $\lambda_1 = \dots = \lambda_\ell = 0$. This is simply because each $\lambda_k$ is an actual entry in the output vector $\Phi(\lambda_1, \dots, \lambda_\ell)$.

Furthermore, we can explicitly state the inverse map $\Phi^{-1}$. For any $x \in \Sol(A, 0)$, we extract the entries corresponding to the free variables:
\begin{equation*}
    \Phi^{-1}(x) = \begin{pmatrix} x_{i_1} \\ \vdots \\ x_{i_\ell} \end{pmatrix}
\end{equation*}
This completes the proof.
\end{proof}

\begin{corollary}[Structure of Non-Homogeneous Solutions]
% cor:17.d.2
\label{cor:non_homogeneous_solutions}
Let $b \in K^m$, and $A \in \M_{m \times n}(K)$. Then:
\begin{enumerate}
    \item $\Sol(A, b) := \{x \in K^n \mid A \cdot x = b\} \neq \emptyset$ if and only if $b \in \ColS(A)$.
    \item If $\Sol(A, b) \neq \emptyset$ and $y \in \Sol(A, b)$, then:
    \begin{equation*}
        \Sol(A, b) = y + \Sol(A, 0) := \{y + x \mid x \in \Sol(A, 0)\}
    \end{equation*}
\end{enumerate}
\end{corollary}

% Extractor: Opus 5
\begin{exercise}[Proof of \cref{cor:non_homogeneous_solutions}]
\label{exc:non_homogeneous_solutions}
The notes leave this proof as an exercise, with the following hints.
% Extractor: Gemini 3.1 Pro
\begin{enumerate}[label=\textbf{(\alph*)}]
    \item Write the matrix $A$ in terms of its column vectors $v_1, \dots, v_n$:
    \[
        A = \begin{pmatrix} | & & | \\ v_1 & \dots & v_n \\ | & & | \end{pmatrix}
    \]
    For any vector $x = \begin{pmatrix} x_1 \\ \vdots \\ x_n \end{pmatrix} \in K^n$, the matrix-vector product is a linear combination of the columns:
    \[
        A \cdot x = x_1 \begin{pmatrix} | \\ v_1 \\ | \end{pmatrix} + \dots + x_n \begin{pmatrix} | \\ v_n \\ | \end{pmatrix}
    \]

    \item Let $y \in \Sol(A, b)$. First, assume $x \in \Sol(A, 0)$. Then $A \cdot x = 0$ and $A \cdot y = b$. Adding these together yields:
    \[
        A(y + x) = A \cdot y + A \cdot x = b + 0 = b
    \]
    Conversely, if $z \in \Sol(A, b)$, meaning $A \cdot z = b$, then define $x := z - y$. We obviously have $z = y + x$, and checking the multiplication gives:
    \[
        A \cdot x = A \cdot (z - y) = A \cdot z - A \cdot y = b - b = 0
    \]
    This shows that any solution $z$ can be written as $y + x$ for some $x \in \Sol(A, 0)$.
\end{enumerate}
\exinfo{This exercise is taken from Prof.~Biran's handwritten lecture notes.}
\exsol{sol:non_homogeneous_solutions}
\end{exercise}

\begin{proposition}[Rank Criterion for Solvability]
% prop:17.d.3
\label{prop:rank_solvability_criterion}
Let $A \in \M_{m \times n}(K)$ and $b \in K^m$. The following statements are equivalent:
\begin{enumerate}
    \item $\rank(A \mid b) = \rank(A)$.
    \item The system of equations $Ax = b$ has a solution.
\end{enumerate}
\end{proposition}

\begin{proof}
Write $A$ in terms of its columns,
\[
    A = \begin{pmatrix} | & & | \\ v_1 & \dots & v_n \\ | & & | \end{pmatrix}.
\]
We have the following chain of equivalences:
\begin{align*}
    \rank(A \mid b) = \rank(A) &\iff \rankcol(A \mid b) = \rankcol(A) \\
    &\iff \dim \Sp(v_1, \dots, v_n, b) = \dim \Sp(v_1, \dots, v_n) \\
    &\iff b \in \Sp(v_1, \dots, v_n) \\
    &\iff b \in \ColS(A) \\
    &\iff \Sol(A, b) \neq \emptyset
\end{align*}
The middle step deserves a word: $\Sp(v_1, \dots, v_n)$ is contained in $\Sp(v_1, \dots, v_n, b)$ whatever $b$ is, and two nested subspaces of equal finite dimension coincide, so the two spans have the same dimension if and only if they are equal, which happens if and only if $b$ already lies in the smaller one. The last step holds directly by \cref{cor:non_homogeneous_solutions}, thus finishing the proof.
\end{proof}

\begin{proposition}[Criterion for a Unique Solution]
% prop:17.d.4
\label{prop:unique_solution_rank_criterion}
Let $A \in \M_{m \times n}(K)$ and $b \in K^m$. The following statements are equivalent:
\begin{enumerate}
    \item $\rank(A) = \rank(A \mid b) = n$.
    \item The system $Ax = b$ has a unique solution.
\end{enumerate}
\end{proposition}

% Extractor: Opus 5
\begin{exercise}[Proof of \cref{prop:unique_solution_rank_criterion}]
\label{exc:unique_solution_rank_criterion}
The notes leave this proof as an exercise, with the following hint.
% Extractor: Gemini 3.1 Pro
Write
\[
    A = \begin{pmatrix} | & & | \\ v_1 & \dots & v_n \\ | & & | \end{pmatrix}.
\]
Note that:
\begin{equation*}
    \rank(A) = n \iff \dim \Sp(v_1, \dots, v_n) = n
\end{equation*}
Since there are $n$ column vectors, this means that $v_1, \dots, v_n$ are linearly independent, and therefore, they form a basis for $\ColS(A)$.
\exinfo{This exercise is taken from Prof.~Biran's handwritten lecture notes.}
\exsol{sol:unique_solution_rank_criterion}
\end{exercise}

% Creator: Opus 5
\subsection{Solutions to the Exercises}
\label{sec:solutions_more_matrices_lineq}

\begin{exercisesolution}[Proof of \cref{cor:non_homogeneous_solutions}]
\label{sol:non_homogeneous_solutions}
\begin{enumerate}[label=\textbf{(\alph*)}]
    \item As the hint records, for $x = (x_1, \dots, x_n)\transp \in K^n$ the product $A \cdot x$ is the linear combination $x_1 v_1 + \dots + x_n v_n$ of the columns of $A$. Hence
    \[
    \Sol(A, b) \neq \emptyset \iff \exists\, x \in K^n : x_1 v_1 + \dots + x_n v_n = b \iff b \in \Sp(v_1, \dots, v_n) = \ColS(A),
    \]
    the last equality being the definition of the column space.

    \item Let $y \in \Sol(A, b)$. If $x \in \Sol(A, 0)$, then $A(y + x) = Ay + Ax = b + 0 = b$, so $y + x \in \Sol(A, b)$, which gives the inclusion $y + \Sol(A, 0) \subseteq \Sol(A, b)$. Conversely, if $z \in \Sol(A, b)$, put $x := z - y$; then $z = y + x$ and
    \[
    Ax = A(z - y) = Az - Ay = b - b = 0,
    \]
    so $x \in \Sol(A, 0)$ and $z \in y + \Sol(A, 0)$. The two inclusions give the asserted equality. \qedhere
\end{enumerate}
\end{exercisesolution}

\begin{exercisesolution}[Proof of \cref{prop:unique_solution_rank_criterion}]
\label{sol:unique_solution_rank_criterion}
Suppose first that $\rank(A) = \rank(A \mid b) = n$. Since the two ranks agree, \cref{prop:rank_solvability_criterion} gives a solution $y \in \Sol(A, b)$. Since $\rank(A) = n$, the computation at the start of this section gives
\[
\dim \Sol(A, 0) = n - \rank(A) = 0,
\]
and therefore, $\Sol(A, 0) = \{0\}$. By \cref{cor:non_homogeneous_solutions} \textbf{(b)}, $\Sol(A, b) = y + \{0\} = \{y\}$, so the solution is unique.

Conversely, suppose $Ax = b$ has a unique solution $y$. In particular, a solution exists, so $\rank(A \mid b) = \rank(A)$ by \cref{prop:rank_solvability_criterion}. If now $x \in \Sol(A, 0)$, then $y + x$ is again a solution of $Ax = b$ by \cref{cor:non_homogeneous_solutions} \textbf{(b)}, and uniqueness forces $y + x = y$, that is, $x = 0$. So, $\Sol(A, 0) = \{0\}$, and therefore
\[
0 = \dim \Sol(A, 0) = n - \rank(A),
\]
which gives $\rank(A) = n$. In the language of the hint: the columns $v_1, \dots, v_n$ are then linearly independent and form a basis of $\ColS(A)$, which is exactly the statement that the coefficients in $b = x_1 v_1 + \dots + x_n v_n$ are uniquely determined.
\end{exercisesolution}

\begin{ainote}
This section is complete in the transcript: all four numbered statements, in the
notes' order, with the notes' proofs and with Prof. Biran's hints for the two he
leaves as exercises. The staircase picture that the notes draw by hand is
rendered as \cref{fig:rref_pivot_staircase}.

The two proofs the notes mark \qt{Exc.} stood in the text as
\verb|\begin{proof}| blocks reading \qt{This is left as an exercise}, followed
by the hints. They are now exercises carrying those hints, and both are solved
above. The second one needs a little more than the hint gives: the uniqueness
criterion follows by combining \cref{prop:rank_solvability_criterion} with the
dimension count $\dim \Sol(A, 0) = n - \rank(A)$ from the opening of the
section.

% Creator: Opus 5
A later pass looking only for gaps in rigour found nothing wrong here. What it
found was the same pattern as in \cref{sec:matrix_multiplication}: an object
defined in running prose with a numbered statement about it standing directly
underneath. The map $\Phi$, which sends a choice of free variables to the
solution it determines, now has a definition of its own, unnumbered so that
17.d.1--17.d.4 are untouched, and \cref{prop:solution_map_isomorphism} points
at it.

Five steps were taken silently and now say why they hold: the dimension of the
solution space rests on the rank theorem together with
$\dim \Img(T_A) = \rank(A)$; the pivot count is
\cref{rem:rank_notation_and_properties}; the word \qt{clearly} in front of
$\Sol(A, 0) = \Sol(A', 0)$ covers \cref{thm:ero_preserve_solutions}, that
row-equivalent systems have the same solutions; the passage from a trivial
kernel to an isomorphism is \cref{cor:equal_dim_iso} \textbf{(c)} followed by
\cref{lem:bijective_is_iso}; and the middle link of the chain of equivalences
in \cref{prop:rank_solvability_criterion} needs the observation that two nested
subspaces of equal finite dimension coincide, without which equal dimensions of
the two spans would not put $b$ in the smaller one.

Finally, the augmented matrix is written $(A \mid b)$ in the chapter on systems
of linear equations, where it is introduced, and in the solutions at the end of
this section, but was written $(A|b)$ in the two propositions. The three
occurrences are unified.
\end{ainote}

